% SPDX-FileCopyrightText: 2026 Bart van de Lint
% SPDX-License-Identifier: AGPL-3.0-or-later
\section{Model inputs}
\label{sec:inputs}

\Cref{tab:inputs} lists every input with its base value, its range where data give one, its unit and one line from its source to its value; where \cref{app:derivations} derives an input, its row points there. A range comes from a measured spread, from sources that disagree or a published uncertainty, or for a goal from what is known today; an input that rests on our own estimate has one value, marked as an estimate, and no range, since a range we invented would make its uncertainty look known. \Cref{tab:development} gives the 15 inputs that the development study of \cref{sec:res_inputs} moves, 18 rows of \cref{tab:inputs}, since it moves the time out of service of three rigs and the two fuel prices each as one input. The ship classes are in \cref{tab:classes}, with each value's source in the published dataset \citep{windshipopt_data}, and the grids, solver settings and constants in \cref{tab:numerics}. The published regressions (Holtrop--Mennen, SNNM, Kijima, ShipCLEAN, the DNV rudder area and the 1969 Tonnage Convention) keep their published coefficients.

\paragraph{Inputs outside the development study} The inputs that \cref{tab:development} leaves out are out for one of five reasons. Those that rest on our estimate have no range to move, among them the price per area or mass of every rig, the canopy's cost, the market factor and the kite's aspect ratio and strong-wind limits (\cref{sec:res_inputs} names those that would matter most if wrong). Most belong to the ship or its setting rather than to a rig's developer: the hull, propeller and engine coefficients, the ship's dimensions, port costs, the yard's share for installing a rig, the discount rate and the smallest ship the carbon price reaches. The tether is bought rope, so its working stress, price and life are the rope maker's. Some are moved through an input that stands in for them: a rig's maintenance through its time out of service, the kite's hardware price through its deck system's mass, and the flight drag through the glide ratio it is calibrated to. The last barely limit any rig of the study: every kite but the Kamsarmax's needs more wind to fly a pattern than the top of the launch wind's range, and that one flies from only \SI{0.1}{m/s} below it (\cref{sec:res_forces}); and wing sails and rotors, built for the ship's \given{25} years, change their yearly cost far less over \given{20 to 30} years of life than canopy life changes a kite's.

{\footnotesize
\setlength{\LTcapwidth}{\textwidth}
\setlength{\extrarowheight}{1.5pt}
\InputIfFileExists{sections/input_tables}{}{\tbd{Input tables pending: run the input\_tables step.}}

\begin{longtable}{@{}>{\raggedright\arraybackslash}p{4.2cm}>{\raggedright\arraybackslash}p{4.8cm}>{\raggedright\arraybackslash}p{6.6cm}@{}}
\caption{Grids, solver settings and constants. Corridor spacings are those of the results, which the draft run shares on a sixth of the departures.}
\label{tab:numerics}\\
\toprule
Setting & Value & Use \\
\midrule
\endfirsthead
\multicolumn{3}{@{}l}{\cref{tab:numerics} (continued)}\\
\toprule
Setting & Value & Use \\
\midrule
\endhead
\bottomrule
\endfoot
\multicolumn{3}{@{}l}{\textit{Performance table}}\\
Speed through the water & 49 speeds evenly from \SI{4}{kn} to the speed at $F_n=0.28$, the same step down to \SI{1.5}{m/s} & thrust and fuel power tables \\
True wind at \SI{10}{m} & 0 to \SI{30}{m/s}, steps of \SI{1}{m/s} & stronger wind is held at \SI{30}{m/s} \\
True wind angle & 0 to \SI{180}{\degree}, steps of \SI{7.5}{\degree} & thrust and fuel power tables \\
Peak period & 3 to \SI{18}{s}, steps of \SI{1}{s} & table of $R_\mathrm{aw}/H_s^2$ \\
Wave heading & 0 to \SI{180}{\degree}, steps of \SI{15}{\degree} & table of $R_\mathrm{aw}/H_s^2$ \\
Wave frequencies & 0.1 to \SI{6}{rad/s}, 240 points & spectral integral of $R_\mathrm{aw}$ \\
Wing trim & $C_L$ 0 to $C_{L,\max}$ attached, angle of attack 25 to \SI{90}{\degree} stalled, solved continuously & best trim \\
Rotor trim & spin ratio 0 to 4 or the surface-speed limit, solved continuously & best trim \\
\multicolumn{3}{@{}l}{\textit{Kite pattern table}}\\
Apparent wind & 3 to \SI{9}{m/s} in steps of \SI{1}{m/s}, then 10.5, 12 and on to \SI{30}{m/s} in steps of \SI{3}{m/s} & pattern search \\
Course off apparent downwind & 0 to \SI{180}{\degree}, steps of \SI{15}{\degree} & pattern search \\
Side-force weights & 0, 0.25, 0.5, 1, 1.5, 2, 3, read linearly between & kite trim, solved continuously \\
Wind reference elevation & \SI{25}{\degree} on the tether & height of the apparent wind of the table \\
Lowest angle of attack & angle where $C_L$ reaches 0.05 & lower bound of the search \\
Points per pattern & 48 & time average of the tether force \\
Pattern search & Nelder--Mead, 200 iterations, three starts and the best pattern of the next stronger wind speed, or the three best of a grid of \num{3360} patterns where none of those flies & kite patterns \\
\multicolumn{3}{@{}l}{\textit{Corridors and routing}}\\
Ocean stage spacing & \SI{45}{nmi}, at least ten stages over a route's open stretches & corridor stages \\
Coastal stage spacing & \SI{30}{nmi} & fixed paths from Rotterdam to Ushant, through the Florida Strait and into each port, no lateral freedom \\
Lateral node spacing & \SI{11.25}{nmi} & nodes on rows square to the great circle of each open stretch \\
Corridor half width & \SI{2400}{nmi} & reach of the nodes to either side of the track, within \ang{60} of latitude \\
Lateral step & up to 5 rows per stage, \SI{51}{\degree} off the corridor's line & edges between stages \\
Planning grid & every stage and row for \SI{240}{nmi} ahead of the ship; beyond, every fourth row to \SI{480}{nmi} and half as many again each time the distance ahead doubles, on stages far enough apart for an edge to reach the next row kept, weather every \SI{6}{h} & each plan \\
Land check along edges & every \SI{20}{nmi} & removes edges over land \\
Coastal clearance & \SI{45}{nmi}, half a wind cell, on rings at 22.5 and \SI{45}{nmi} of 16 bearings & no open-sea node nearer land \\
Port approach & \SI{90}{nmi}, one wind cell, from a port or the end of a coastal path & open-sea nodes kept near land \\
Island disks & island radius plus about \SI{20}{nmi} & land mask \\
Time bin & \SI{3}{h} & state of the dynamic programme \\
Weather along an edge & at most every \SI{3}{h}, \SI{6}{h} beyond the first \SI{240}{nmi} of a plan & fuel and speed of an edge \\
States searched & within 3\,\% of a backward estimate of the cost to go at \SI{24}{h} knots, then 6\,\%, then all, if a pass ends above its limit or finds no path; costs at most 0.01 points of cost saving against searching all states, 0.35 for the kite & dynamic programme \\
Longest passage & 50 days & time bins \\
Speeds & every eighth table speed, 7 from \SI{4}{kn} to the ship's maximum, every sixteenth in the backward estimate, and on a long edge those between that bring its arrivals within one time bin of each other & edge choices \\
Final pass & all table speeds, \SI{0.25}{h} bins, \SI{1}{h} on the forecast plan weighing the wind error, within \SI{6}{h} of the path found & passage reported \\
Live sailing & speed and kite chosen every \SI{1}{h} over every fourth table speed and the plan's, on the forecast result & planner on forecasts \\
Re-planning interval & first route node \SI{24}{h} after the last plan; 12 and \SI{6}{h} in the re-planning study & planner on forecasts \\
Plan margin on the kite's limits & \SI{20}{\%} tighter wind limits and more launch wind in a sea, tested & planner on forecasts \\
Forecasts & GEFSv12 reforecast control run and four perturbed members, wind and waves, issued daily to 16 days, mean bias against ERA5 removed per grid point and lead, blended with ERA5's 1991--2016 monthly climatology by a weight fitted per grid point and lead & planner on forecasts \\
Operational forecasts & GEFS operational control run, wind, issued four times a day, mean bias against ERA5 removed and blended with the climatology in the same way; waves from the daily reforecast & re-planning study \\
Departures & every 14 days from 1 January 2017 to 100 days before the end of 2019; in all runs the one nearest the 15th of each calendar month (12) & weather sample \\
\multicolumn{3}{@{}l}{\textit{Constants}}\\
Gravitational acceleration & \SI{9.81}{m/s^2} & Froude number and waves \\
Earth radius & \SI{6371}{km} & great-circle distances \\
Nautical mile & \SI{1852}{m} & distances and speeds in knots \\
Hours in a year & 8760 & annual costs \\
\end{longtable}}

\paragraph{Rig sales} \stale{classes}{The rig sales of \cref{sec:industry} (\cref{tab:market_sales}) count every ship of a class as fitting, at its best size, the rig that wins on a route, so a sixth of the class's ships count for each route the rig wins. The low end of a range counts the classes and routes a rig wins alone, the high end adds those where it is level with another; the best rig fits whichever saves most cost, from the classes and routes one rig wins alone to all where some rig clearly saves, both by the rule of \cref{sec:res_robustness}. Sales are the price of fitting those ships once at the prices of \cref{sec:costs}, without spare canopies. A kite costs more than a rotor, but a ship carries one kite against several rotors, so on e-ammonia the rotors sell about twice as much as the kite although the kite suits a larger part of shipping (\cref{fig:fleet_share}).}

\stale{classes}{By ship type (\cref{tab:market_sales_types}): \tbd{The rotors' sales are mostly to bulk carriers and then tankers, on both fuels. The kite's are spread over bulk carriers, container ships and crude oil tankers, so a kite maker would enter on container ships, the one type where the kite leads or is level with the rotors on both fuels.}}

\begin{table*}[tbp]
\centering
\footnotesize
\caption{\stale{classes}{} Ships and rig sales of \cref{tab:market_sales} by ship type for the kite and the rotors, the rigs that win some class and route, on each fuel, from the classes and routes the rig wins alone to those where it is level with another rig too, by the rule of \cref{sec:res_robustness}. Estimates on the same fleet data and assumptions; a dash marks no sales.}
\label{tab:market_sales_types}
\InputIfFileExists{sections/market_sales_type_table}{}{\tbd{Market sales type table pending: run the market\_table step.}}
\end{table*}
